\documentclass[11pt]{article}
\usepackage[margin=1in]{geometry}
\usepackage{amsmath,amssymb,amsfonts,mathtools}
\usepackage{array,booktabs,enumitem,graphicx,xcolor,hyperref}
\usepackage[T1]{fontenc}
\usepackage[utf8]{inputenc}
\title{Bi-Hamiltonian Sturcture of Super KP Hierarchy}
\author{Source-backed arXiv sample}
\date{}
\begin{document}
\maketitle
\begin{abstract}
We obtain the bi-Hamiltonian structure of the super KP hierarchy based on the even super KP operator \$\textbackslash{}Lambda = \textbackslash{}theta\textasciicircum{}\{2\} + \textbackslash{}sum\textasciicircum{}\{\textbackslash{}infty\}\_\{i=-2\}U\_\{i\} \textbackslash{}theta\textasciicircum{}\{-i-1\}\$, as a supersymmetric extension of the ordinary KP bi-Hamiltonian structure. It is expected to give rise to a universal super \$W\$-algebra incorporating all known extended superconformal \$W\_\{N\}\$ algebras by reduction. We also construct the super BKP hierarchy by imposing a set of anti-self-dual constraints on the super KP hierarchy.
\end{abstract}
\section{Source Metadata}
This sample is based on arXiv record \texttt{hep-th/9109009}. Authors: Feng Yu.
\section{Extracted Prose 1}
We obtain the bi-Hamiltonian structure of the super KP hierarchy based on the even super KP operator math expression , as a supersymmetric extension of the ordinary KP bi-Hamiltonian structure. It is expected to give rise to a universal super math expression -algebra incorporating all known extended superconformal math expression algebras by reduction. We also construct the super BKP hierarchy by imposing a set of anti-self-dual constraints on the super KP hierarchy.
\section{Extracted Prose 2}
The integrable nonlinear (partial) differential KdV [1] and KP [2] equations are known to possess bi-Hamiltonian structures. This was first conjectured by Adler [3] and then proved by Gelfand and Dickey [4] for the KdV hierarchy, via the famous Gelfand-Dickey bracket. For the KP hierarchy, which is a set of multi-time evolution equations in Lax form
\end{document}
